\chapter{Introdução} 
\label{cap:1}

\indent A análise estrutural é o ramo da mecânica dos sólidos que estuda o comportamento de estruturas quando a elas são aplicados esforços. Em muitos casos práticos de engenharia esses esforços apresentam variações dos esforços pouco significativas em relação ao tempo, permitindo que a estrutura seja analisada estaticamente. Todavia, os esforços que variam em relação ao tempo ou que mobilizem as forças inerciais da estrutura devem ter uma análise mais apurada considerando os efeitos dinâmicos.\\
\indent Em estruturas civis, a análise dos efeitos dinâmicos é de suma importância seja quanto ao conforto dos usuários quanto a estabilidade das estruturas. Quando forças externas atuam com uma frequência seja próxima a uma das frequências naturais da estrutura, ocorre o efeito de ressonância que produz excessivos deslocamentos e possíveis falhas (\citet{mechanicalvibrations2017}). Mesmo garantindo a estabilidade, a verificação do conforto dos usuários apresenta um papel crucial na análise da estrutura e através da análise dinâmica pode-se verificar a necessidade de enrijecer a estrutura ou adicionar sistemas de amortecimento para atendimento desse critério.\\
\indent Diferente da análise estática, que envolve a solução de um sistema algébrico de equações, a obtenção dos resultados para a análise dinâmica é uma tarefa muito mais árdua, mesmo para modelos simples. Na engenharia civil era comum a análise dinâmica ser evitada, substituindo-a por uma análise estática com coeficientes de majoração dinâmicos. 
Mas, com o passar do tempo, as estruturas têm tornado-se cada vez mais esbeltas e, portanto, mais suscetíveis a vibrações, tornando 
imprescindível a realização da análise dinâmica de estruturas. (\citet{modelagemestruturasacomistas2017}).\\
\indent O processo envolvido na análise estrutural se inicia com a elaboração de um modelo mecânico representativo da estrutura. Os modelos podem ser classificados como analíticos, estatísticos, ou mistos, e essa distinção ocorre através da forma como as equações são obtidas, em sendo que os modelos analíticos são modelados utilizando um princípio físico enquanto modelos empíricos, ou estatísticos, são obtidos com ensaios (\citet{licoesemmecanicaestruturasdinamica2022a}). Na análise dinâmica, um modelo estatístico é construído, por exemplo, através de um ensaio de túnel de vento enquanto um modelo analítico é formulado através dos princípios da mecânica do contínuo.\\
\indent Solucionar de forma exata um modelo analítico, com exceção daqueles com geometrias simples ou condições de contorno específicas, é um processo trabalhoso e em muitos casos impossível, e para isso é comum o emprego de métodos numéricos para obtenção de uma solução aproximada. Um dos métodos mais amplamente utilizados é o Método dos Elementos Finitos (MEF) que particiona um domínio complexo em subdomínios que podem ser aproximados. A aplicação do método dos MEF na análise de estruturas, potencializados por um aumento do poder computacional, permitiu a elaboração de modelos estruturais muito mais complexos se comparados aos que eram utilizados até 40 anos atrás (\citet{mechanicalvibrations2017}), porém, vale ressaltar, que a resposta obtida -- assim como por outras técnicas numéricas -- é aproximada. \\
\indent O MEF, apesar de já consagrado na comunidade científica e técnica, apresenta restrições relacionadas a geometria da malha, como os elementos possuírem numero de nós e padrão geométrico fixo, a depender da família de elementos, e elementos muito distorcidos resultam em erros na solução numérica. O desfecho são dificuldades de modelagem e solução de problemas onde existe a reconfiguração da malha, resultando em elementos cuja distorção pode afetar a precisão da solução. Buscando solucionar o problema de malhas foi desenvolvido o Método dos Elementos Virtuais (MEV), que possibilita a construção de elementos de formato arbitrário mantendo a mesma estrutura do método de Galerkin utilizado pelo MEF. Como destacado em \citet{engineeringperspectivevirtualelementmethoditsinterplaystandardfiniteelementmethod2019a}, o MEV possui Entre suas vantagens está a capacidade de possuir arestas com ângulos entre si que podem ser de 180° ou convexos.  A etapa de pré-processamento do modelo é mais rápida considerando que não existe preocupação com a conformidade da malha, ou seja, uma malha que não possua elementos com grandes distorções, principalmente ao levar em conta que a elaboração de uma malha conforme demanda considerável quantidade de tempo.\\
\indent Em trabalhos anteriores desenvolvidos nos programas de pós-graduação da UFPR, foram realizadas análises dinâmicas de problemas em estado plano comparando o MEF com o Método dos Elementos Finitos Generalizado (MEFG) (\citet{torii2012analise},\citet{cittadin2020enriquecimento} e \citet{custodio2022metodo}). Este trabalho busca ser pioneiro na utilização do MEV dentro do programa de Pós-graduação em Métodos Numéricos em Engenharia e comparar a acurácia dos resultados obtidos pelo método com os trabalhos previamente desenvolvidos por outros alunos.

\section{Objetivos}

\indent O objetivo deste trabalho é investigar o comportamento do MEV na análise de vibração 
livre em problemas de estado plano de tensões e deformações.

\section{Objetivos específicos}

\indent Pa alcançar o objetivo geral deste trabalho, pretende-se:

\begin{enumerate}
    \item Desenvolver a formulação do MEV em soluções lineares e de alta ordem.
    \item Implementar o MEV para análise de vibração livre e validar os resultados com modelos de referência.
    \item Comparar os resultados obtidos pelo MEV com as soluções do MEF e MEFG.
\end{enumerate}

\section{Estrutura do trabalho}

\indent Este trabalho será dividido em 6 capítulos organizados da seguinte maneira: no segundo capítulo é feita a revisão da literatura do MEV destacando os principais evoluções do método. O capítulo 3 aborda a análise dinâmica de estruturas, os estados planos de tensões e deformações e apresenta o desenvolvimento da formulação variacional. NO capítulo 4 é explicada a formulação e as propriedades do MEV assim como a forma de computar as matrizes de massa e de rigidez. No capítulo 5 são abordados três problemas de vibração: uma chapa engastada trapezoidal com furo circular no centro no estado plano de tensões, uma barragem de terra no estado plano de deformações e uma parede estrutural de 4 pavimentos no estado plano de tensões. As conclusões obtidas são contempladas no último capítulo deste trabalho.











